\chapter{Benchmark and Report Artifact Schema}
\label{app:result-schema}

\section{Benchmark JSON structure}

The benchmark harness writes a JSON object with two top-level keys:
\begin{lstlisting}[language=Python,caption={Conceptual benchmark JSON schema.}]
{
  "schema_version": 1,
  "metadata": { ... },
  "results": [
    { ... one benchmark row ... },
    ...
  ]
}
\end{lstlisting}

\paragraph{Metadata block.}
The \texttt{metadata} object records run-level information such as timestamp, command line, git state, software versions, CUDA device properties, and benchmark CLI configuration.

\paragraph{Results block.}
Each entry in \texttt{results} corresponds to one method-shape-dtype-causality configuration. Rows preserve both success and failure states. This is important: a report may summarize only successful rows, but the raw artifact should retain OOMs, unavailable methods, and invalid configurations.

\section{Representative row fields}

\begin{longtable}{@{}p{0.30\textwidth} p{0.16\textwidth} p{0.46\textwidth}@{}}
\caption{Representative fields in one benchmark result row.}
\label{tab:result-row-fields}\\
\toprule
Field & Example type & Meaning \\
\midrule
\endfirsthead
\toprule
Field & Example type & Meaning \\
\midrule
\endhead
\texttt{status} & string & \texttt{ok}, \texttt{error}, \texttt{oom}, \texttt{unavailable}, or \texttt{invalid\_configuration} \\
\texttt{method} & string & \texttt{cuda}, \texttt{sdpa}, \texttt{reference}, or \texttt{auto} \\
\texttt{dtype} & string & input tensor dtype name \\
\texttt{query\_length}, \texttt{key\_length}, \texttt{head\_dim} & integers & problem shape parameters \\
\texttt{causal} & boolean & causal or non-causal mode \\
\texttt{mode} & string & \texttt{forward} or \texttt{forward-backward} \\
\texttt{latency\_ms\_median} & float or null & median timing sample for successful rows \\
\texttt{dense\_equivalent\_tflops} & float or null & simplified FLOP-rate estimate \\
\texttt{incremental\_peak\_memory\_mib} & float or null & incremental allocator peak in MiB \\
\texttt{max\_abs\_error} & float or null & maximum absolute forward error versus the reference \\
\texttt{reference\_status} & string & whether reference validation ran, was skipped, or failed \\
\texttt{error} & string or null & exception text for unsuccessful rows \\
\bottomrule
\end{longtable}

\section{Rendered report snippets}

The helper script \texttt{benchmarks/render\_report\_artifacts.py} derives the following optional files from one benchmark JSON artifact:
\begin{itemize}
  \item \texttt{report/generated/environment.tex}
  \item \texttt{report/generated/correctness\_summary.tex}
  \item \texttt{report/generated/performance\_summary.tex}
  \item \texttt{report/generated/memory\_summary.tex}
  \item \texttt{report/generated/profiling\_summary.tex}
\end{itemize}

Each generated file should:
\begin{enumerate}
  \item state the source benchmark artifact path;
  \item contain only derived information from that artifact;
  \item remain safe to delete and regenerate;
  \item never become the sole location of a result that is absent from the raw JSON.
\end{enumerate}

\section{Expected workflow}

\begin{lstlisting}[language=bash,caption={Result-to-report workflow.}]
python -m benchmarks.bench_attention \
  --device cuda --methods cuda sdpa reference \
  --seq-lens 128 512 1024 --head-dims 64 --dtypes float16 --causal both \
  --output-dir results/runs --require-cuda

python -m benchmarks.render_report_artifacts \
  --benchmark-json results/runs/benchmark_<run-id>.json \
  --output-dir report/generated

cd report
latexmk -pdf main.tex
\end{lstlisting}

\section{Why a schema appendix belongs in the report}

This appendix may seem unusual for a course-style report, but it is valuable for a systems portfolio project. The point is to make the evidence pipeline inspectable. Another reader should be able to trace a table in Chapter~\ref{chap:results} back to a concrete JSON artifact and then back to the exact benchmark command that produced it.
